\begin{pastexam}[label={px:ccu-114-swa-07}]{中正資工所甲組 114 學年度 第 7 題}
（複選）

(5\%) (Select \underline{multiple} correct answers) Which of the following statements is (are) true?
\begin{enumerate}[label=(\alph*)]
\item Insertion sort performs in-place sorting operations.
\item Finding the longest simple path from an unweighted graph can be solved by divided and conquer
strategy.
\item Recursion tree is a full binary tree.
\item Quicksort is not an asymptotically optimal algorithm.
\item An algorithm with fewer lines of codes has lower time complexity than the one with more lines of
codes.
\end{enumerate}
\end{pastexam}

\begin{pastsoln}
答案：(a)(d)。

本題混合了排序（資結書第~14~章）、動態規劃的最佳子結構（第~18.1.4~節）和遞迴樹（第~2~章）。

\begin{itemize}
\item (a) 對。 插入排序只在原陣列裡把元素往後挪，額外只用 $O(1)$ 的變數（資結書第~1.3~節），是原地
排序。
\item (b) 錯。 分治要把問題切成獨立的子問題，各自的最佳解合起來就是整體的最佳解。 最長簡單路徑沒有
這種最佳子結構（資結書第~18.1.4~節）：$q \to r \to t$ 是 $q$ 到 $t$ 的最長簡單路徑時，$q \to r$ 這段
不一定是 $q$ 到 $r$ 的最長簡單路徑；兩段各自最長的路徑合起來，可能重複經過同一個頂點，就不是簡單路徑了。
例如四個點的環 $q - r - t - s - q$：$q$ 到 $r$ 最長的簡單路徑是 $q, s, t, r$，$r$ 到 $t$ 最長的是
$r, q, s, t$，接起來重複了。 這個問題是 NP 難的（判斷是否有 $n - 1$ 條邊的簡單路徑，就是漢米爾頓路徑
問題），一般認為沒有多項式時間的演算法。
\item (c) 錯。 遞迴樹每個節點的子節點數是遞迴式裡子問題的個數：$T(n) = 3T(n/4) + n^2$ 的遞迴樹每個節點
有 $3$ 個子節點，$T(n) = T(n-1) + n$ 只有 $1$ 個。 不一定是二元樹，更不一定是 full binary tree。
\item (d) 對。 「漸近最佳」指最壞情況的時間達到比較排序的下界 $\Omega(n \log n)$；合併排序、堆積排序
是，快速排序最壞情況 $\Theta(n^2)$（資結書第~14.5~節），所以不是。 它的平均情況 $\Theta(n \log n)$
是最佳的，但說一個演算法漸近最佳，看的是最壞情況。
\item (e) 錯。 程式行數和時間複雜度無關。 例如氣泡排序的程式比合併排序短，卻是 $\Theta(n^2)$ 對
$\Theta(n \log n)$。
\end{itemize}
\end{pastsoln}
